\documentclass[11pt]{article}
\usepackage[margin=1in]{geometry}
\usepackage{amsmath,amssymb,amsthm,mathtools}
\usepackage{bm}
\usepackage{enumitem}
\usepackage{hyperref}
\usepackage{xcolor}
\usepackage{listings} % Include listings package

% Setup the listings package for Python syntax highlighting
\lstset{language=Python,
  basicstyle=\ttfamily\small,
  keywordstyle=\color{blue},
  commentstyle=\color{olive},
  stringstyle=\color{red},
  showstringspaces=false,
  identifierstyle=\color{black}}


\hypersetup{colorlinks=true,linkcolor=blue,citecolor=blue,urlcolor=blue}

\title{LQR Control of a Vibrating Square Membrane\\\large Project Handout}
\author{}
\date{}

\newcommand{\R}{\mathbb{R}}
\newcommand{\Om}{\Omega}
\newcommand{\dd}{\,d}
\newcommand{\p}{\partial}
\newcommand{\ip}[2]{\left\langle #1,#2\right\rangle}
\newtheorem*{remark}{Remark}

\begin{document}
\maketitle

\section*{Project summary}
We consider a square membrane with small transverse displacement. The 
membrane is fixed on its boundary and is actuated by pressing on a 
single localized spot. The goal is to design a linear quadratic regulator 
(LQR) that drives the membrane to rest.

This project begins with a partial differential equation and reduces it 
to a finite-dimensional control system by expanding in Laplacian 
eigenfunctions. The finite-dimensional truncation is then treated with 
standard LQR machinery.

\section{PDE model}
Let
\[
\Omega = (0,1)\times(0,1)
\]
be the unit square. Let $u(x,y,t)$ denote the vertical displacement of 
the membrane. For small displacements, $u$ satisfies the controlled 
wave equation
\begin{equation}
\label{eq:wave}
 u_{tt}(x,y,t) = c^2 \Delta u(x,y,t) + b(t)\,\psi(x,y), 
 \qquad (x,y)\in\Omega,\ t>0,
\end{equation}
with fixed boundary conditions
\begin{equation}
\label{eq:bc}
 u(x,y,t)=0,\qquad (x,y)\in \partial\Omega,
\end{equation}
and initial data
\begin{equation}
\label{eq:ic}
 u(x,y,0)=u_0(x,y),\qquad u_t(x,y,0)=v_0(x,y).
\end{equation}
Here $c>0$ is the wave speed, $b(t)\in\R$ is the scalar control input, 
and $\psi(x,y)$ is the actuator profile.

The ideal model of pressing at one point $(x_0,y_0)\in\Omega$ is 
obtained by taking
\[
\psi(x,y)=\delta(x-x_0)\delta(y-y_0),
\]
but this is singular. For numerical work, it is better to use a 
localized patch actuator, for example, a narrow Gaussian or a bump 
function centered at $(x_0,y_0)$.

\begin{remark}
For this project, the point-actuator model is useful for the 
mathematics, while a sharply localized smooth profile is better for 
computation.
\end{remark}

\section{Eigenfunction expansion on the square}
The Dirichlet Laplacian eigenfunctions on the unit square are
\begin{equation}
\label{eq:eigs}
\phi_{mn}(x,y)=2\sin(m\pi x)\sin(n\pi y),\qquad m,n\ge 1,
\end{equation}
with eigenvalues
\begin{equation}
\label{eq:eigvals}
\lambda_{mn}=\pi^2(m^2+n^2).
\end{equation}
These eigenfunctions form an orthonormal basis of $L^2(\Omega)$.

We expand
\begin{equation}
\label{eq:expansion}
 u(x,y,t)=\sum_{m=1}^{\infty}\sum_{n=1}^{\infty} q_{mn}(t)\,\phi_{mn}(x,y).
\end{equation}
Substituting \eqref{eq:expansion} into \eqref{eq:wave} and projecting 
onto mode $(m,n)$ gives
\begin{equation}
\label{eq:modalode}
 \ddot q_{mn}(t) + c^2\lambda_{mn} q_{mn}(t) = \beta_{mn}\, b(t),
\end{equation}
where
\begin{equation}
\label{eq:beta-general}
\beta_{mn}=\int_0^1\!\int_0^1 \psi(x,y)\phi_{mn}(x,y)\,\dd x\,\dd y.
\end{equation}
If we use the ideal point actuator at $(x_0,y_0)$, then
\begin{equation}
\label{eq:beta-point}
\beta_{mn}=\phi_{mn}(x_0,y_0)=2\sin(m\pi x_0)\sin(n\pi y_0).
\end{equation}

Thus every retained mode behaves like a harmonic oscillator driven by 
the same scalar input $b(t)$, but with mode-dependent gain $\beta_{mn}$.

\section{Finite-dimensional truncation}
Choose an integer $M\ge 1$ and retain modes with
\[
1\le m,n\le M.
\]
There are then $N=M^2$ retained modes. Flatten the double index 
$(m,n)$ into a single index $k=1,\dots,N$, and write
\[
q(t)=\begin{bmatrix} q_1(t) \\ \vdots \\ q_N(t) \end{bmatrix},
\qquad
p(t)=\dot q(t).
\]
Define
\[
\omega_k^2=c^2\lambda_k,
\qquad
\beta=\begin{bmatrix}\beta_1 \\ \vdots \\ \beta_N\end{bmatrix},
\qquad
\Omega_N = \operatorname{diag}(\omega_1^2,\dots,\omega_N^2).
\]
Then the truncated system becomes
\begin{equation}
\label{eq:firstorder}
\dot x = A x + B b,
\qquad
x=\begin{bmatrix} q \\ p \end{bmatrix}\in\R^{2N},
\end{equation}
with
\begin{equation}
\label{eq:AB}
A=
\begin{bmatrix}
0 & I \\
-\Omega_N & 0
\end{bmatrix},
\qquad
B=
\begin{bmatrix}
0 \\
\beta
\end{bmatrix}.
\end{equation}

\section{Energy and LQR cost}
The truncated mechanical energy is
\begin{equation}
\label{eq:energy}
E(t)=\frac12\sum_{k=1}^N \left(p_k(t)^2 + \omega_k^2 q_k(t)^2\right).
\end{equation}
A natural infinite-horizon quadratic cost is
\begin{equation}
\label{eq:cost}
J(b)=\int_0^{\infty} \left( x(t)^T Q x(t) + R\,b(t)^2 \right)\dd t,
\end{equation}
where $Q\succeq 0$ and $R>0$.

A physically meaningful choice is
\begin{equation}
\label{eq:Qchoice}
Q=
\begin{bmatrix}
\alpha\,\Omega_N & 0 \\
0 & \beta_v I
\end{bmatrix},
\qquad \alpha>0,\ \beta_v>0,
\end{equation}
so that displacement is penalized according to modal potential energy 
and velocity is penalized according to kinetic energy.

\section{LQR feedback law}
For the finite-dimensional system \eqref{eq:firstorder}, the 
infinite-horizon LQR feedback has the form
\begin{equation}
\label{eq:lqr-law}
 b(t)=-Kx(t),
\end{equation}
where
\begin{equation}
\label{eq:gain}
K=R^{-1}B^T P,
\end{equation}
and $P$ is the symmetric positive semidefinite solution of the 
algebraic Riccati equation
\begin{equation}
\label{eq:ARE}
A^T P + P A - P B R^{-1} B^T P + Q = 0.
\end{equation}
The closed-loop dynamics are
\begin{equation}
\label{eq:closedloop}
\dot x = (A-BK)x.
\end{equation}
If the truncated pair $(A,B)$ is stabilizable and $(Q^{1/2},A)$ is 
detectable, then the closed-loop matrix $A-BK$ is stable and the 
modal amplitudes decay to zero.

\subsection{Note on Riccati equation and ``direct method''}
Riccati equation will not be used in this project. You will use
the ``direct method'' which only relies upon the ability to
find eigenvalues. However, you need to be familiar with Riccati
equation, as it is often the best way to solve small examples,
such as ones that often appear on exams.

\section{Why actuator location matters}
The actuator enters mode $(m,n)$ through the factor
\[
\beta_{mn}=2\sin(m\pi x_0)\sin(n\pi y_0)
\]
for the point-actuator model.

If $\beta_{mn}=0$, then that mode is not directly actuated. This is 
especially important for symmetric actuator locations. At the center of 
the square,
\[
(x_0,y_0)=\left(\frac12,\frac12\right),
\]
we get
\[
\beta_{mn}=2\sin\left(\frac{m\pi}{2}\right)\sin\left(\frac{n\pi}{2}\right),
\]
which vanishes whenever $m$ or $n$ is even. Thus a centered actuator 
misses many modes.

For this reason, a good default choice is an off-center actuator such 
as
\[
(x_0,y_0)=(0.37,0.61).
\]
That breaks symmetry and typically produces nonzero coupling to the 
first many modes.

\section{A concrete model for the students}
A convenient starter model is:
\begin{itemize}[leftmargin=1.5em]
\item membrane domain $\Omega=(0,1)^2$,
\item wave speed $c=1$,
\item actuator location $(x_0,y_0)=(0.37,0.61)$,
\item truncation size $M=6$ or $M=8$,
\item point-actuator couplings from \eqref{eq:beta-point},
\item LQR weights from \eqref{eq:Qchoice} with, say, $\alpha=1$, 
$\beta_v=1$, $R=10^{-2}$ or $10^{-1}$.
\end{itemize}

\section{Repository guide}
The starter repository includes a Python path.
In the Python code, the main driver scripts are:
\begin{itemize}[leftmargin=1.5em]
\item \texttt{src/python/run\_demo.py} for a first closed-loop 
simulation,
\item \texttt{src/python/scan\_actuator.py} for comparing actuator 
locations,
\item \texttt{src/python/modal\_lqr.py} for model construction, 
LQR design, and membrane reconstruction.
\end{itemize}
The repository also includes a task list with milestones. Students 
are encouraged to keep their code modular so that changing the 
truncation size, actuator location, or cost weights requires only a 
few edits.

\section{Project tasks}
Possible project goals include:
\begin{enumerate}[leftmargin=1.5em]
\item Derive the modal equations from the PDE.
\item Build the truncated matrices $A$ and $B$.
\item Compute the LQR gain $K$.
\item Simulate open-loop and closed-loop trajectories for selected 
initial conditions.
\item Reconstruct the membrane shape from the modal coefficients.
\item Investigate the role of actuator location.
\item Compare a point actuator with a smooth localized actuator 
patch.
\end{enumerate}

\section{Possible research questions}
A mathematically meaningful research question is:
\begin{quote}
How does the location of a single localized actuator affect 
controllability and LQR stabilization of a vibrating square membrane?
\end{quote}
More specific directions:
\begin{itemize}[leftmargin=1.5em]
\item What modes are lost when the actuator is placed at the center?
\item How does the choice of $R$ affect the tradeoff between fast 
stabilization and control effort?
\item How large should the truncation be before the observed behavior 
becomes qualitatively stable?
\item Does a regularized actuator patch improve numerical robustness 
compared with the ideal point-actuator model?
\end{itemize}

\section{Suggested deliverables}
A final submission should include at least the following core items:
\begin{itemize}[leftmargin=1.5em]
\item a derivation of the finite-dimensional model,
\item code that constructs the LQR controller,
\item plots of modal amplitudes and control input,
\item snapshots or an animation of the membrane displacement,
\item a short discussion of actuator placement and uncontrollable modes.
\end{itemize}

\subsection{IMPORTANT: Forbidden Python packages} 
The final submission is not allowed to use these imports:
\begin{lstlisting}[language=Python]
from scipy.integrate import solve_ivp
from scipy.linalg import solve_continuous_are
\end{lstlisting}

You must implement your own version of continuous LQR solver and
IVP solver. Your LQR solver must be based on ``direct method''
discussed in class and the ``numbook''. Thus, you need to implement these
two functions:

\begin{enumerate}
\item \texttt{solve\_continuous\_are}
\item \texttt{solve\_ivp}
\end{enumerate}

Make sure that the signatures of these functions and values returned
correspond to their usage in the file \texttt{modal\_lqr.py}.

The Gradescope Autograder will call your implementation
through functions:

\begin{enumerate}
  \item \texttt{simulate\_closed\_loop}
  \item \texttt{simulate\_open\_loop}
\end{enumerate}

You must put your implementation in the file
\texttt{student.py} in folder \texttt{src/python}.

\subsection{Functions not to be modified}
The autograder will call the function \texttt{build\_model}.
You should not modify this function, with one exception: if
you are implementing the optional extension described below.


\section*{Optional extension}
If desired, one may add viscous damping to the PDE:
\[u_{tt}=c^2\Delta u - \gamma u_t + b(t)\psi(x,y),\qquad \gamma>0.\]

Then the modal equations become
\[
\ddot q_k + \gamma \dot q_k + \omega_k^2 q_k = \beta_k b(t).
\]
This is not necessary for the project, but it can make the numerical 
simulations more forgiving.

\subsection{How to handle the extension, so that the autograder notices it?}
The autograder hook is the function \texttt{build\_model}. Add the
parameter \(\gamma\) to your model if you implement this optional extension.
Autograder will call your \texttt{build_model}.





\end{document}
